% Source: https://artofproblemsolving.com/community/c3h2511845_mmc_and_mdc_of_such_numbers_where_one_of_them_is_0_zero_it_is_possible_
% Stable identity: aops-topic-2511845-q01
\textbf{(a)} Every integer $d$ divides $0$, since $0=d\cdot0$. So the common divisors of $A$ and $0$ are exactly the divisors of $A$, and the greatest of them is $A$: $\gcd(A,0)=\boxed{A}$.

\textbf{(b)} The multiples of $0$ are the numbers $0\cdot k=0$, so the only multiple of $0$ is $0$ itself. Hence the only common multiple of $A$ and $0$ is $0$, which is not positive, and there is no least positive common multiple.

\textbf{(c)} With $\operatorname{lcm}(A,0)=0$:
\begin{itemize}
\item $\gcd(A,0)\cdot\operatorname{lcm}(A,0)=A\cdot0=0=A\cdot0$, so the identity still holds;
\item the only common multiple of $A$ and $0$ is $0$, which is a multiple of $0$, so every common multiple is a multiple of the lcm.
\end{itemize}
So the convention $\operatorname{lcm}(A,0)=0$ keeps the usual properties: $0$ is the ``least'' common multiple in the sense of divisibility, though not in size.
